% Новые переменные, которые могут использоваться во всём проекте
% ГОСТ 7.0.11-2011
% 9.2 Оформление текста автореферата диссертации
% 9.2.1 Общая характеристика работы включает в себя следующие основные структурные
% элементы:
% актуальность темы исследования;
\newcommand{\actualityTXT}{Актуальность темы исследования.}
% степень ее разработанности;
\newcommand{\progressTXT}{Степень разработанности темы исследования.}
% цели и задачи;
\newcommand{\aimTXT}{Целью диссертационной работы}
\newcommand{\tasksTXT}{задачи}
% Объект и субъект исследования
\newcommand{\researchObjectTXT}{Объектом исследования}
\newcommand{\researchSubjectTXT}{Предметом исследования}
% научную новизну;
\newcommand{\noveltyTXT}{Научной новизной}
% научную новизну;
\newcommand{\theoreticalValueTXT}{Теоретическая значимость}
\newcommand{\practicalValueTXT}{Практическая значимость}
% теоретическую и практическую значимость работы;
%\newcommand{\influenceTXT}{Теоретическая и практическая значимость}
% или чаще используют просто
\newcommand{\influenceTXT}{Практическая значимость}
% методологию и методы исследования;
\newcommand{\methodsTXT}{Методология и методы исследования.}
% положения, выносимые на защиту;
\newcommand{\defpositionsTXT}{Положения, выносимые на защиту:}
% положения, выносимые на защиту;
\newcommand{\specialityRelationTXT}{Область исследования.}
% степень достоверности и апробацию результатов.
\newcommand{\reliabilityTXT}{Степень достоверности}
\newcommand{\probationTXT}{Апробация работы.}

\newcommand{\contributionTXT}{Личный вклад автора.}
\newcommand{\publicationsTXT}{Публикации.}


%%% Заголовки библиографии:

% для автореферата:
\newcommand{\bibtitleauthor}{\centering{СПИСОК РАБОТ, ОПУБЛИКОВАННЫХ АВТОРОМ ПО ТЕМЕ ДИССЕРТАЦИИ}}% (ГОСТ Р 7.0.11-2011, 9.3)

% для стиля библиографии `\insertbiblioauthorgrouped`
\newcommand{\bibtitleauthorvak}{\centerline{В изданиях из списка ВАК РФ}}
\newcommand{\bibtitleauthorscopus}{\centerline{В изданиях, входящих в международную базу цитирования Scopus}}
\newcommand{\bibtitleauthorwos}{\centerline{В изданиях, входящих в международную базу цитирования Web of Science}}
\newcommand{\bibtitleauthorother}{\centerline{В прочих изданиях}}
\newcommand{\bibtitleauthorconf}{\centerline{В сборниках трудов конференций}}
\newcommand{\bibtitleauthorpatent}{\centerline{Зарегистрированные патенты}}
\newcommand{\bibtitleauthorprogram}{\centerline{Зарегистрированные программы для ЭВМ}}

% для стиля библиографии `\insertbiblioauthorimportant`:
\newcommand{\bibtitleauthorimportant}{Наиболее значимые \protect\MakeLowercase\bibtitleauthor}

% для списка литературы в диссертации и списка чужих работ в автореферате:
\newcommand{\bibtitlefull}{Список литературы} % (ГОСТ Р 7.0.11-2011, 4)
